We modeled the kinematics of the T-Hex kit and created a simulation package for testing locomotion controllers. We implemented three different controllers, along with a posture adaptation pipeline that gives the nominal base height for locomotion. We integrated the posture adaptation pipeline with the open-loop kinematic gait and tested it in flat terrain and found success in navigating confined spaces. However, the open-loop kinematic gait was expectedly poor at walking in rough terrain. The DRL-based controller was successful in learning to walk in rough terrain, albeit with some limitations in the random blocks rough terrain. The MIT Cheetah 3-like controller did not perform well, likely due to the assumptions made in the dynamic model used for control. 

Our next steps include integrating the posture adaptation pipeline with the DRL-based controller, building upon the MIT Cheetah 3-like controller to adapt it for our robot, and testing the DRL-based controller and the MIT Cheetah 3-like controller in hardware. To integrate the posture adaptation pipeline with the DRL-based controller, we plan on training three separate policies for three different nominal base heights, and then use the posture adaptation pipeline to determine which nominal base height to use at any given time. To adapt the MIT Cheetah 3-like controller for our robot, we plan on looking into more accurate dynamic models that take into account the mass of the legs, as well as more robust approaches to foothold selection. Furthermore, as the controller generates torque commands, we also plan to adapt the controller to convert torque commands into position commands \cite{zepedaTheoreticalConversionTorque2024,delpreteImplementingTorqueControl2016,khatibTorquepositionTransformerTask2008} so that it can be tested on hardware.